%% ============================================================
%%  CHAPTER 2 — CANONICAL SCAFFOLD
%%  Capacity Utilization in the Center and the Periphery:
%%  Old Debates and New Estimates
%%
%%  Status   : SCAFFOLD — prose stubs marked [DRAFT PROSE]
%%  Equations: LOCKED from RRR_Accumulation_Framework.md
%%             and capital_accumulation_dynamics_v5.md
%%  Spine    : RRR_Accumulation_Framework.md
%%  Outline  : Ch2_Outline_v1_DRAFT.md (§2.1–§2.11)
%%  Date     : 2026-03-29
%% ============================================================

\documentclass[12pt,letterpaper]{article}

%% -- Packages ------------------------------------------------
\usepackage{amsmath, amssymb, amsthm}
\usepackage{mathtools}
\usepackage{booktabs}
\usepackage{array}
\usepackage{graphicx}
\usepackage{setspace}
\usepackage[margin=1in]{geometry}
\usepackage[authoryear,round]{natbib}
\usepackage{hyperref}
\usepackage{xcolor}

%% -- Custom commands (locked notation) ----------------------
\newcommand{\E}{\mathbb{E}}
\newcommand{\R}{\mathbb{R}}
\newcommand{\pd}[2]{\frac{\partial #1}{\partial #2}}
\newcommand{\pdd}[2]{\frac{\partial^2 #1}{\partial #2^2}}

%% Structural objects (Ch2 locked)
\newcommand{\mut}{\mu_t}           % capacity utilization
\newcommand{\Yp}{Y^p_t}            % productive capacity
\newcommand{\thetaL}{\theta(\Lambda)}  % transformation elasticity
\newcommand{\pit}{\pi_t}           % profit share
\newcommand{\iotat}{\iota_t}       % gross accumulation rate
\newcommand{\betat}{\beta_t}       % recapitalization rate
\newcommand{\Bt}{B_t}              % capital productivity at normal capacity
\newcommand{\phit}{\phi_t}         % investment share

%% Prose stub marker (REMOVE IN FINAL)
\newcommand{\draftnote}[1]{%
  \textcolor{red}{\textbf{[DRAFT PROSE: #1]}}}

%% -- Theorem environments -----------------------------------
\newtheorem{proposition}{Proposition}[section]
\newtheorem{definition}{Definition}[section]

%% -- Document -----------------------------------------------
\begin{document}

\begin{titlepage}
  \vspace*{2cm}
  \begin{center}
    {\Large \textbf{Chapter 2}}\\[0.5em]
    {\large Capacity Utilization in the Center and the Periphery:\\
    Old Debates and New Estimates}\\[2em]
    {\normalsize Dissertation Chapter — Canonical Scaffold}\\[0.5em]
    {\normalsize Version: 2026-03-29}
  \end{center}
  \vspace{1cm}
  \begin{abstract}
    \draftnote{Abstract: 150--200 words. State: (1)~the structural deficiency
    of existing CU estimates and why it matters for crisis theory; (2)~the
    RRR/CVAR identification strategy and what it recovers; (3)~the behavioral
    recapitalization law and its nested closures; (4)~the Tavares peripheral
    twist as the chapter's distinctive contribution; (5)~the three downstream
    objects inherited by Ch3.}
  \end{abstract}
\end{titlepage}

\doublespacing

%% ===========================================================
\section{Introduction}
\label{sec:intro}
%% ===========================================================
%
% Arc: crisis theory → CU as the missing structural link →
%      two-stage empirical machine → center/periphery comparison
%
% Deliverables flagged: (1) mu_hat, (2) B_hat + theta_hat, (3) behavioral law

Capacity utilization occupies a peculiar position in macroeconomic analysis:
universally invoked, rarely identified. The profit rate decomposes as
$r_t = \mut \cdot \Bt \cdot \pit$ --- utilization scales normal profitability
--- so the path of $\mut$ mediates every link between accumulation dynamics and
crisis. Yet the capacity utilization series that enter empirical work are
typically constructed by statistical filters that impose balanced growth by
assumption, setting the capacity transformation elasticity $\theta$ to unity
and suppressing the regime-dependent content through which structural
disequilibrium operates. Chapter~1 demonstrated this deficiency: existing CU
measures --- from the CBO output gap to the Banco Central de Chile's structural
balance framework --- recover a mean-reverting residual by construction, not a
structurally identified object. This chapter asks a different question:
\emph{what does capacity utilization do in the economy, and how does its role
differ between center and periphery?}

The answer proceeds through a two-stage empirical machine applied identically
to the United States and Chile. Stage~A recovers capacity utilization as a
latent confinement object. The structural primitives are accounting identities:
output decomposes into productive capacity and utilization
($y_t = y^p_t + \ln\mut$), productive capacity is linked to the capital stock
through the transformation elasticity ($y^p_t = \thetaL\, k_t$), and the
confinement identity $\ln\mut = y_t - \thetaL\, k_t$ follows immediately.
The parameter $\thetaL$ is an upstream regime parameter pinned by the
institutional space --- not a behavioral choice and not derivable from any
first-order condition. The empirical specification embeds distribution through
the interaction term $\pit k_t$, yielding the long-run capacity
relation~\eqref{eq:cv} with endogenous elasticity
$\theta_t = \beta_1 + \beta_2\pit$. If the residual is stationary, the
recovered $\widehat{\mut}$ is the structurally identified utilization series ---
and with it come the derived objects $\hat{B}_t$ (capital productivity at normal
capacity) and $\hat{y}^p_t$ (productive capacity) that Stage~B requires as
inputs.

Stage~B estimates a behavioral accumulation law. The accounting identity
$\iotat = \betat\, r_t$ decomposes the gross accumulation rate into the
recapitalization rate and the profit rate; the behavioral question is whether
$\betat$ responds to distribution, demand, and the structural transformation of
capital productivity. The behavioral accumulation
law~\eqref{eq:accumulation_law} nests three closures --- Cambridge,
Bhaduri--Marglin, and full structural --- testable via Wald restrictions on the
coefficients $b_j = 1 + \eta_j$. The deviation from unity is the behavioral
content.

The same machine is run on both countries, and the structural comparison is
meaningful because the measurement apparatus is identical. The key asymmetry
that emerges is not a measurement artefact but a structural finding: in the
periphery, the recapitalization response to capital productivity is constrained
by the availability of foreign exchange. The Tavares channel --- the chapter's
distinctive theoretical contribution --- localizes this constraint in the
recapitalization function rather than in the output equation. Peripheral
accumulation requires imported machinery; the capital-goods sector is
structurally incomplete; and $\betat$ collapses when the foreign currency needed
to purchase imported capital goods is unavailable
\citep{tavares1985}. The coefficient $b_3$ --- the recapitalization response to
capital productivity --- is the empirical locus of this constraint, and the
comparison $b_3^{US}$ versus $b_3^{CL}$ gives the Tavares channel testable
content.

The remainder of the chapter is organized as follows.
Section~\ref{sec:crisis} establishes the crisis-theoretic motivation.
Section~\ref{sec:primitives} develops the accounting scaffold.
Section~\ref{sec:distribution} introduces the distributive micro-foundation.
Section~\ref{sec:bridge} bridges the Ch1 focal data pair to Ch2's state vector.
Section~\ref{sec:stageA} estimates the structural identification (Stage A).
Section~\ref{sec:accounting} derives the accounting law of capital accumulation.
Section~\ref{sec:stageB} estimates the behavioral recapitalization law (Stage B).
Section~\ref{sec:tavares} formalizes the peripheral twist.
Section~\ref{sec:comparison} synthesizes the center--periphery comparison.
Section~\ref{sec:conclusions} concludes.


%% ===========================================================
\section{Crisis Theory: Stagnation Tendencies, Cumulative Disequilibrium,
  and the Role of Capacity Utilization}
\label{sec:crisis}
%% ===========================================================

Chapter~1 established that existing capacity utilization measures are
structurally deficient: they impose balanced growth by construction and
suppress the regime-dependent content of the capital--output relationship.
This section establishes why that deficiency matters. Capacity utilization is
the mechanism through which accumulation dynamics trigger crisis --- and
without a structurally identified $\mut$, the crisis mechanism cannot be read
empirically.


%%------------------------------------------------------------
\subsection{Stagnation and Crisis: A Typology}
\label{sec:crisis_typology}
%%------------------------------------------------------------

The empirical analysis of capacity utilization in this chapter is motivated by
a prior question: what role does utilization play in the dynamics of capitalist
crisis? Before that question can be answered, the varieties of crisis must be
distinguished. The stagnation typology developed by
\citet{vidal2014,vidal2019} and \citet{basu2019} organizes the structural
coordinates within which crises unfold and clarifies where capacity utilization
enters the causal chain.

The typology is structured as a two-by-two matrix along two axes
(Table~\ref{tab:typology}). The rows distinguish the source of disruption:
demand-side mechanisms, in which the realization of surplus value is
interrupted, and financial-side mechanisms, in which the circuits of credit and
asset valuation that mediate accumulation break down. The columns distinguish
the distributional configuration: an excess of surplus value, in which profits
are abundant relative to the economy's capacity to absorb them productively,
and a deficit of surplus value, in which profits are squeezed by rising labor
strength or deteriorating productivity conditions.

Each cell identifies a distinct crisis tendency. Under-consumption arises when
the surplus is generated but cannot be realized through effective demand.
Profit squeeze describes the configuration in which rising labor strength
compresses the profit share, eroding the incentive to accumulate. Financial
fragility emerges when abundant surplus is channeled into speculative asset
positions rather than productive investment, generating Minsky-type
instability. The falling rate of profit --- the classical Marxian tendency ---
operates when capital deepening outpaces productivity gains and the profit rate
declines secularly.

These four tendencies are not mutually exclusive; concrete historical episodes
typically combine elements of more than one cell. What matters for the present
chapter is the distinction between partial and structural crisis. A partial
crisis is an interruption resolvable by partial reconfiguration of the
institutional space~$\Lambda$ --- a policy adjustment, a shift in the terms of
class compromise, a financial restructuring. A structural crisis requires
drastic restructuring of the accumulation regime itself: a change in the
technology of production, the organization of labor markets, the international
division of labor, or the financial architecture.

The typology locates the structural coordinates within which crises develop,
but it does not specify the mechanism that triggers the transition from
accumulation to crisis. That trigger --- the cumulative disequilibrium in the
path of capacity utilization --- is the subject of the next subsection.

\begin{table}[h]
  \centering
  \caption{Stagnation Tendencies: A Typology}
  \label{tab:typology}
  \begin{tabular}{lcc}
    \toprule
     & Excess of Surplus Value & Deficit of Surplus Value \\
    \midrule
    Demand side     & Under-consumption / realization crisis
                    & Profit squeeze (rising labor strength) \\
    Financial side  & Financial fragility
                    & Falling (normal) rate of profit \\
    \bottomrule
  \end{tabular}
  \bigskip

  \noindent{\small \textit{Notes}: Following \citet{vidal2014,vidal2019}.
  Each cell denotes a qualitatively distinct stagnation tendency.
  Partial crisis $\equiv$ interruption resolvable by partial reconfiguration
  of the institutional space $\Lambda$.
  Structural crisis $\equiv$ deep interruption requiring drastic restructuring.}
\end{table}


%%------------------------------------------------------------
\subsection{Okishio's Crisis Trigger and Cumulative Disequilibrium}
\label{sec:okishio}
%%------------------------------------------------------------

The stagnation typology of~\S\ref{sec:crisis_typology} identifies the
structural coordinates within which crises unfold, but it does not specify a
trigger. What sets a crisis in motion? \citeauthor{okishio1961}'s
(\citeyear{okishio1961}) contribution provides the answer that organizes this
chapter: decentralized capitalist investment decisions produce a secular
cumulative disequilibrium in capacity utilization, and it is the \emph{path}
of~$\mut$ --- not a static threshold --- that triggers the transition from
partial adjustment to structural crisis.

The core of Okishio's argument is that the mismatch between effective demand
and productive capacity is not self-correcting. Individual capitalists,
responding to current profitability signals, expand capacity in ways that are
locally rational but globally destabilizing. When demand is strong and
utilization is high, investment responds --- but the resulting capacity
additions do not arrive instantaneously, and when they do, the aggregate demand
conditions that motivated them may no longer prevail. The consequence is a
secular tendency for $\mut$ to trace an unstable path: episodes of high
utilization attract investment that eventually depresses utilization, while
episodes of low utilization discourage investment and allow demand to recover,
but each cycle leaves a residue of structural imbalance. Utilization rates
trend toward a reversal of normal capacity but overshoot, producing stagnation
episodes of increasing severity.

The critical insight is that the crisis trigger is the path of $\mut$, not a
level. There is no fixed utilization rate at which crisis erupts. Rather, the
cumulative trajectory --- the sequence of overshoots and undershoots, the
widening amplitude of cyclical swings, the secular drift of the utilization
path relative to the capacity manifold --- is itself the destabilizing
mechanism. This is what distinguishes Okishio's framework from static
equilibrium models in which utilization gravitates toward a given normal rate:
the disequilibrium is cumulative, and its accumulation is the crisis.

The channel through which the utilization path transmits into crisis runs
through profitability. The \citeauthor{weisskopf1979}
(\citeyear{weisskopf1979}) decomposition of the profit rate provides the
accounting link:
\begin{equation}
  r_t = \mut \cdot \Bt \cdot \pit,
  \label{eq:weisskopf}
\end{equation}
where $r_t \equiv \Pi_t/K_t$ is the profit rate, $\mut \equiv Y_t/\Yp$ is
capacity utilization, $\Bt \equiv \Yp/K_t$ is capital productivity at normal
capacity, and $\pit \equiv \Pi_t/Y_t$ is the profit share.
Profitability is demand-led through $\mut$: utilization scales normal
profitability $\Bt\pit$ up or down, so that the entire profit rate fluctuates
with the utilization path. When $\mut$ falls persistently --- when the
cumulative disequilibrium that Okishio describes tilts toward chronic
underutilization --- profitability erodes even if $\pit$ remains stable and
$\Bt$ does not deteriorate. The crisis is transmitted from the demand side to
the profitability side through the utilization channel, and from profitability
to accumulation through the investment response.

This transmission mechanism makes the identification of~$\mut$ a precondition
for any empirical analysis of the crisis--accumulation nexus. To
read~\eqref{eq:weisskopf} as an empirical statement rather than a tautology,
one needs a structurally identified utilization series --- one that is not
imposed by a statistical filter but recovered from the long-run reproduction
structure of the economy. The existing CU estimates examined in Chapter~1
impose balanced growth by construction ($\theta = 1$), suppressing precisely
the regime-dependent content through which cumulative disequilibrium operates.
The structural primitives of~\S\ref{sec:primitives} provide the alternative: a
confinement identity in which $\mut$ is recovered as the deviation of realized
output from the productive-capacity manifold under unbalanced growth. The next
subsection traces the link from profitability to accumulation that completes
the crisis mechanism.


%%------------------------------------------------------------
\subsection{From Profitability to Accumulation: The Missing Link}
\label{sec:missing_link}
%%------------------------------------------------------------

The preceding subsection established that the path of~$\mut$ transmits
cumulative disequilibrium into profitability through the Weisskopf
decomposition~\eqref{eq:weisskopf}. But the crisis mechanism is incomplete
without a further link: profitability must translate into accumulation for the
disequilibrium to propagate. The literature on the tendency of the profit rate
to fall \citep{basu2012,basu2013} has documented the long-run trajectory
of~$r_t$ in the US and other advanced economies, yet the causal channel from
profitability to crisis runs not through the profit rate itself but through the
investment response it elicits. As \citet{basu2017} and
\citet{ibarra2023,ibarra2024} argue, the rate of profit matters for crisis
dynamics because it drives the rate of investment --- and it is the rate of
investment, not the rate of profit, that determines whether the economy
reproduces or contracts its productive capacity.

This chapter formalizes that link through the behavioral accumulation law
developed in~\S\ref{sec:accounting}--\ref{sec:stageB}. The gross accumulation
rate $\iotat = I_t/K_t$ decomposes as $\iotat = \betat\, r_t$, where $\betat$
is the recapitalization rate. The behavioral question is whether $\betat$
responds to the structural conditions summarized by $\pit$, $\mut$, and $\Bt$.
The estimation of that response --- carried out in Stage~B of the empirical
machine --- requires structurally identified objects: the utilization series
$\widehat{\mut}$ and the capital-productivity series $\hat{B}_t$ recovered from
Stage~A.

The mechanism differs structurally between center and periphery. In the US, the
recapitalization rate responds to the three profitability channels with
relatively few binding external constraints. In Chile, the foreign-exchange
constraint mediates the conversion of surplus into investment goods, so that the
recapitalization response to capital productivity --- the $b_3$ coefficient ---
is structurally dampened. The Tavares channel (\S\ref{sec:tavares}) provides
the theoretical mechanism for that dampening.

To make the crisis-trigger hypothesis testable, $\mut$ must be identified
structurally --- not imposed by a statistical filter. Chapter~1 showed why
existing measures fail; the next section formalizes the alternative.


%% ===========================================================
\section{Structural Primitives: Production, Capacity,
  and Unbalanced Growth}
\label{sec:primitives}
%% ===========================================================

This section establishes the accounting scaffold on which all subsequent
identification and estimation rests. Every relation derived here is a
definition or an identity; no behavioral equation is imposed
until~\S\ref{sec:accounting}. The reader should treat each step as a
manipulation of national-accounts magnitudes, not as a claim about agent
optimization.

We adopt a Leontief production structure in which both capital and labor bind
at full capacity utilization:
\begin{equation}
  Y_t = \min\!\left(B_t K_t,\; A_t L_t\right),
  \label{eq:leontief}
\end{equation}
with both constraints binding at $\mut = 1$. Output decomposes into
productive capacity and utilization:
\begin{align}
  Y_t &= \mut \cdot \Yp, \label{eq:decomp}\\
  y_t &= y^p_t + \ln\mut, \label{eq:decomp_log}
\end{align}
where lowercase denotes natural logarithms of the corresponding uppercase
level, $y_t \equiv \ln Y_t$, $k_t \equiv \ln K_t$.
Capital productivity at normal capacity is
\begin{equation}
  \Bt \equiv \frac{\Yp}{K_t},
  \qquad
  \ln \Bt = y^p_t - k_t.
  \label{eq:caprod}
\end{equation}

The long-run transformation of capital stock into productive capacity is
governed by the \emph{unbalanced growth closure}:
\begin{equation}
  y^p_t = \thetaL \cdot k_t,
  \label{eq:ugclosure}
\end{equation}
where $\thetaL$ is the capacity transformation elasticity — a regime
parameter, pinned by the institutional space $\Lambda$ and not derivable
from within-period optimization.
Balanced growth ($\thetaL = 1$) is the knife-edge special case;
the general case is $\thetaL \neq 1$.

\begin{proposition}[Confinement Identity]
\label{prop:confinement}
Given the decomposition~\eqref{eq:decomp_log} and the unbalanced growth
closure~\eqref{eq:ugclosure},
\begin{equation}
  \ln\mut = y_t - \thetaL\, k_t.
  \label{eq:confinement}
\end{equation}
Capacity utilization is identified as the deviation of realized output from
the productive-capacity manifold.
\end{proposition}

The confinement identity is the core identification object of the chapter.
Three features merit emphasis. First, it is pure accounting: no behavioral
assumption has been invoked, only the decomposition of output into capacity and
utilization and the definition of the capacity transformation elasticity.
Second, when $\thetaL \neq 1$ the identity implies persistent structural
disequilibrium in the capital--output ratio --- the economy does not settle on a
balanced-growth path, and the deviation of~$\theta$ from unity measures the
degree of structural imbalance. Under-mechanization ($\theta > 1$) implies that
capital accumulation generates more than proportional capacity expansion;
over-mechanization ($\theta < 1$) implies that the economy accumulates capital
faster than it creates productive capacity, so that the capital--output ratio
rises secularly. Third, the Harrodian benchmark $\theta = 1$ is not the default
but the special case: imposing it amounts to assuming away precisely the
structural content that this chapter seeks to recover.

To take the confinement identity to data, one requires an empirically tractable
specification of~$\theta$. The next section introduces the
distribution-conditioned capacity relation, in which the transformation
elasticity varies with the profit share, and derives the rank-1 identification
equation that makes $\mut$ recoverable from observed time series.


%% ===========================================================
\section{Distribution-Conditioned Capacity: The Interaction Term}
\label{sec:distribution}
%% ===========================================================

The confinement identity derived in~\S\ref{sec:primitives} recovers capacity
utilization as the deviation of realized output from the productive-capacity
manifold, but it leaves $\thetaL$ as an abstract regime parameter. To take the
identity to data, one needs an empirically tractable specification that links
$\theta$ to observable magnitudes. This section introduces the
distribution-conditioned capacity relation, in which the transformation
elasticity varies with the profit share, and derives the rank-1 identification
equation that makes $\mut$ recoverable from observed time series.

Let the exploitation rate $e_t$ map into the profit share as
\begin{equation}
  \pit = \frac{e_t}{1 + e_t},
  \label{eq:profitshare}
\end{equation}
and introduce the interaction term $z_t \equiv \pit \cdot k_t$ as the
empirical proxy for $\theta(e|\Lambda)\cdot k_t$ in the cointegrating space.
The long-run capacity relation then takes the estimable form:
\begin{equation}
  y_t = c + \beta_1 k_t + \beta_2(\pit k_t) + \xi_t,
  \qquad
  \xi_t \equiv \ln\mut,
  \label{eq:cv}
\end{equation}
so that the endogenous long-run transformation elasticity is
\begin{equation}
  \theta_t = \beta_1 + \beta_2\, \pit.
  \label{eq:theta_linear}
\end{equation}

\begin{proposition}[Rank-1 Identification]
\label{prop:rank1}
If $\xi_t = y_t - c - \beta_1 k_t - \beta_2(\pit k_t)$ is stationary, then
\begin{equation}
  \widehat{\ln\mut} = y_t - \hat{c} - \hat{\beta}_1 k_t - \hat{\beta}_2(\pit k_t)
  \label{eq:cv_estimated}
\end{equation}
is the recovered latent utilization object. This is the dissertation's
primary identification problem.
\end{proposition}

Three features of this identification deserve emphasis. First, $\theta$ remains
upstream: the parameters $\beta_1$ and $\beta_2$ are structural parameters of
the cointegrating relation, not behavioral margins subject to agent
optimization. The linearization~\eqref{eq:theta_linear} is the empirical proxy
for the general $\theta(e|\Lambda)$ introduced
in~\S\ref{sec:primitives}; the capacity transformation elasticity is a
property of the institutional regime, not a choice variable. Second, the
interaction term $z_t = \pit k_t$ rotates the cointegrating space --- it is a
constructed object needed to represent $\theta(e)\cdot k$ within the linear
structure of the Johansen VECM. Without it, the cointegrating relation would
impose a constant $\theta = \beta_1$, collapsing the distribution-conditioned
specification back to the regime-invariant case. Third, the Increasing
Intensity of Conflict Schedule (IICS) is a falsifiable hypothesis, not an
imposed condition. The sign of~$\beta_2$ is unrestricted: if $\beta_2 > 0$,
then higher exploitation raises $\theta_t$ and distributional polarization
expands the economy's capacity-generating ability per unit of capital; if
$\beta_2 < 0$, the opposite holds; if $\beta_2 = 0$, distribution does not
condition the transformation elasticity and the specification reduces to the
confinement identity of~\S\ref{sec:primitives} with a constant slope. The data
adjudicate. The next section embeds this identification in the
center--periphery comparison that structures the empirical work of the chapter.


%% ===========================================================
\section{From Chapter~1 to Chapter~2: The Focal Data Pair
  and the State Vector}
\label{sec:bridge}
%% ===========================================================

Chapter~1 established the bivariate output--capital pair as the empirically
validated focal data for US capacity utilization, diagnosed rank-1 confinement
via ARDL bounds tests, and recovered a preliminary transformation elasticity.
This section maps that focal pair into the trivariate cointegrating state
vector of Chapter~2, extends the construction to Chile, and states the
methodological rationale for the parallel design.


%%------------------------------------------------------------
\subsection{US: From Ch1's ARDL to Ch2's Cointegrating Relation}
\label{sec:bridge_us}
%%------------------------------------------------------------

\draftnote{§2.5.1: 1.5 pages. Ch1's bivariate ARDL on $(y_t, k_t)$ with
  VAcorp/KGCcorp recovered $\hat{\theta}^{S1}$ and diagnosed rank-1
  confinement (S2). Ch2 augments with $z_t = \pit k_t$, allowing $\theta_t$
  to vary with distribution. State vector:
  $\mathbf{X}_t^{US} = (y_t,\, k_t,\, \pit k_t)'$.
  Data: BEA/FRED $(y,k)$ from Ch1; NIPA Table~1.12 $(e \to \pi)$.
  Sample: 1947--2019. Positioning: contrast with CBO, Adams--Coe,
  multivariate filter methods that impose $\theta=1$.}

The state vector for the US system is
\begin{equation}
  \mathbf{X}_t^{US} = \begin{pmatrix} y_t \\ k_t \\ \pit k_t \end{pmatrix},
  \label{eq:statevec_us}
\end{equation}
estimated on annual data from 1947 to 2019, with output and capital stock
inherited from the Chapter~1 BEA/FRED pipeline, and the profit share
constructed from NIPA Table~1.12.


%%------------------------------------------------------------
\subsection{Chile: Constructing the Parallel}
\label{sec:bridge_chile}
%%------------------------------------------------------------

The Chilean empirical construction mirrors the US case in identification
strategy while confronting distinct data challenges and a fundamentally
different institutional context for capacity utilization measurement. The data
construction requires splicing across multiple statistical vintages. The
Chilean capital stock series draws on \citet{hofman2000} for the pre-1960
period, the Penn World Table (\texttt{rnna}) for the post-1960 benchmark, and
the \citet{luders2016} and \citet{perez2019} series for interpolation and
cross-validation. The splicing procedure matches growth rates at overlap years
rather than levels, preserving the long-run trend properties that the
cointegration analysis requires. The profit share is constructed from national
accounts data on compensation of employees, spliced with the
\citet{alarco2014} adjusted wage-share series for the pre-1990 period and
cross-checked against the \citet{astorga2023} historical distributional series.
The exploitation rate $e_t$ is then recovered from
$\pit = e_t/(1+e_t)$, and the interaction term
$z_t^{CL} = \pi^{CL}_t k^{CL}_t$ is constructed exactly as in the US case.
The estimation sample runs from 1960 to 2019, constrained at the lower end by
capital-stock availability.

The positioning of this identification against the existing Chilean CU
literature is consequential. Capacity utilization in Chile is not merely an
academic object --- it is a governance variable that feeds directly into the
structural fiscal balance rule administered by the Banco Central de Chile and
the Direcci\'on de Presupuestos. The BCCh framework, formalized in
\citet{fuentes2008} and updated in \citet{figueroa2019} and
\citet{durand2024}, estimates potential output using multivariate filter methods
that impose stationarity of the output gap by construction. The
\citet{magud2011} IMF estimates for Chile adopt a similar approach. In all
these frameworks, the transformation elasticity $\theta$ is implicitly set to
unity: capital and capacity are assumed to grow at the same rate, and the output
gap is a mean-reverting deviation from a balanced-growth trend.

The identification developed in this chapter offers an alternative. By
permitting $\theta \neq 1$ and recovering the distribution-conditioned capacity
slope, it allows the data to reveal whether the Chilean economy operates under
over-mechanization ($\theta < 1$) --- a possibility that the structural-balance
framework rules out by assumption. The key asymmetry between the US and Chilean
cases --- that CU estimates in the US serve primarily as academic benchmarks
while in Chile they govern fiscal policy directly --- is not an artefact of
measurement but a finding that emerges from the institutional embedding of the
concept.

The state vector for the Chilean system mirrors~\eqref{eq:statevec_us}:
\begin{equation}
  \mathbf{X}_t^{CL} = \begin{pmatrix} y^{CL}_t \\ k^{CL}_t \\ \pi^{CL}_t k^{CL}_t \end{pmatrix}.
  \label{eq:statevec_cl}
\end{equation}


%%------------------------------------------------------------
\subsection{The Parallel Construction as Method}
\label{sec:bridge_method}
%%------------------------------------------------------------

Both countries receive the same state vector
$\mathbf{X}_t = (y_t, k_t, \pit k_t)'$, the same Johansen cointegration
procedure, and the same rank-1 identification equation derived
in~\S\ref{sec:distribution}. The structural comparison
in~\S\ref{sec:comparison} is meaningful because the measurement apparatus is
identical: any difference in the estimated~$\theta$, in the recovered
$\widehat{\mut}$, or in the behavioral coefficients $b_1, b_2, b_3$ reflects a
structural difference between the two economies, not a difference in how the
objects were constructed.

This is the methodological core of the parallel design. The US budget-baseline
embedding of CU --- in which utilization estimates inform Congressional Budget
Office projections but carry no direct fiscal consequence --- and the Chilean
structural-balance embedding --- in which CU estimates govern the fiscal
expenditure ceiling --- are not parallel by construction. They are asymmetric
institutional facts. That asymmetry is a finding of the comparison, not an
artefact of the method. The parallel construction isolates the source of
difference: because the identification is the same, the differences that emerge
in~\S\ref{sec:comparison} can be attributed to the structural and institutional
conditions under which accumulation proceeds, not to the econometric choices
that produced the estimates.


%% ===========================================================
\section{Stage A: Structural Identification (RRR/CVAR)}
\label{sec:stageA}
%% ===========================================================

Stage~A estimates the long-run capacity relation derived
in~\S\ref{sec:distribution} and recovers the latent utilization
object~$\widehat{\mut}$. The econometric framework is the
\citet{johansen1988} maximum-likelihood procedure applied to
$\mathbf{X}_t = (y_t, k_t, \pit k_t)'$, with rank determination by trace and
maximum-eigenvalue tests, a restricted constant specification, and
\citet{gregoryhansen1996} endogenous break detection. The estimation setup is
common to both countries; results are reported separately below.


%%------------------------------------------------------------
\subsection{Estimation}
\label{sec:stageA_est}
%%------------------------------------------------------------

The \citet{johansen1988} procedure embeds the long-run capacity relation
within a vector error-correction model (VECM) of the form
\begin{equation}
  \Delta \mathbf{X}_t = \Pi \mathbf{X}_{t-1}
    + \sum_{i=1}^{p-1}\Gamma_i \Delta\mathbf{X}_{t-i}
    + \Phi D_t + \varepsilon_t,
  \label{eq:vecm}
\end{equation}
where $\Pi = \alpha\beta'$ is the reduced-rank matrix that encodes the
cointegrating relations, $\Gamma_i$ are short-run adjustment matrices, and
$D_t$ collects deterministic components. The rank of~$\Pi$ determines the
number of cointegrating vectors. Under the identification
of~\S\ref{sec:distribution}, the expected rank is $r = 1$: a single long-run
relation ties output, capital, and the interaction term together, and the
stationary residual of that relation is $\ln\mut$.

Rank determination proceeds through two complementary test sequences. The
trace test evaluates $H_0\colon r \leq j$ against $H_1\colon r > j$ for
$j = 0, 1, 2$, while the maximum-eigenvalue test evaluates $H_0\colon r = j$
against $H_1\colon r = j+1$. Critical values follow
\citet{johansenjuselius1990}, adjusted for the finite-sample correction of
\citet{reimers1992} given the moderate sample lengths involved (roughly 60--70
annual observations for the US and Chile).

The deterministic specification is the restricted constant model: a drift term
is absorbed into the cointegrating space rather than left unrestricted in the
common trends. This choice reflects the theoretical prior that the long-run
capacity relation admits a non-zero intercept --- the constant~$c$
in~\eqref{eq:cv} --- but that no deterministic linear trend should enter the
cointegrating residual. Intuitively, $\ln\mut$ fluctuates around a mean but
does not trend; the restricted constant ensures that the econometric
specification respects this property.

Structural-break detection employs the \citet{gregoryhansen1996} test for
cointegration with an endogenous break at an unknown date. This is particularly
relevant for Chile, where the 1973 military coup and the 1982 debt crisis
represent candidate regime shifts that may alter the cointegrating parameters.
For the US, the candidate break corresponds to the Fordist/post-Fordist
transition around~1973. The GPY~(2025) impossibility result --- which proves
that the cointegrating vector~$\beta$ cannot be fully nonparametric ---
provides formal justification for maintaining constant cointegrating vectors
within institutional regimes and testing for breaks at regime boundaries rather
than estimating continuously time-varying parameters.


%%------------------------------------------------------------
\subsection{Results: US Baseline}
\label{sec:stageA_us}
%%------------------------------------------------------------

\draftnote{§2.6.2: 2--3 pages (to be filled after estimation).
  Report: $\hat{\beta}_1$, $\hat{\beta}_2$ $\Rightarrow$ $\hat{\theta}_t^{US}$.
  Recovered $\widehat{\ln\mut^{US}}$. Regime classification: does $\hat{\theta}^{US}$
  cross unity? When? Loading coefficients and speed of adjustment.
  Diagnostics: CUSUM, recursive eigenvalues.}


%%------------------------------------------------------------
\subsection{Results: Chile}
\label{sec:stageA_chile}
%%------------------------------------------------------------

\draftnote{§2.6.3: 2--3 pages (to be filled after estimation).
  Same structure as §2.6.2. First structural comparison:
  $\hat{\theta}^{US}$ vs.\ $\hat{\theta}^{CL}$ — the mechanization gap in
  $\theta$-units. Discuss regime breaks: pre-/post-1973, pre-/post-1982.}


%%------------------------------------------------------------
\subsection{Derived Structural Objects}
\label{sec:stageA_derived}
%%------------------------------------------------------------

Once the cointegrating parameters are recovered, the following structural
objects are constructed by definition:

\begin{definition}[Recovered Productive Capacity]
\label{def:yp}
\begin{equation}
  \hat{y}^p_t = y_t - \widehat{\ln\mut}.
  \label{eq:yp_hat}
\end{equation}
\end{definition}

\begin{definition}[Recovered Capital Productivity at Normal Capacity]
\label{def:B}
\begin{equation}
  \widehat{\ln\Bt} = \hat{y}^p_t - k_t,
  \qquad
  \hat{B}_t = \exp\!\left(\hat{y}^p_t - k_t\right).
  \label{eq:B_hat}
\end{equation}
\end{definition}

These are the Stage B inputs: $\widehat{\mut}$ and $\hat{B}_t$ are
passed into the behavioral accumulation law.

The recovered series $\hat{y}^p_t$ is the productive-capacity path implied by
the cointegrating relation --- the output level that would prevail if
$\mut = 1$. The recovered $\hat{B}_t$ measures the efficiency of the capital
stock evaluated at normal utilization, net of cyclical demand fluctuations. Its
trajectory over time is the empirical counterpart of the unbalanced-growth
dynamics discussed in~\S\ref{sec:primitives}: when $\thetaL < 1$, $\hat{B}_t$
declines secularly as capital accumulates faster than productive capacity
expands.

The two-stage architecture is load-bearing at this point. An ARDL investment
equation of the type estimated in Stage~B cannot perform Stage~A's
identification --- it can describe how accumulation responds to utilization, but
it cannot recover what utilization \emph{is}. The utilization
object~$\widehat{\mut}$ and the capital-productivity object~$\hat{B}_t$ are
available cleanly only after the cointegrating relation has been estimated and
the rank-1 identification confirmed. Section~\ref{sec:accounting} derives the
accounting law that decomposes accumulation into its structural components,
and~\S\ref{sec:stageB} specifies the behavioral law that uses these recovered
objects as inputs.


%% ===========================================================
\section{The Accounting Law of Capital Accumulation}
\label{sec:accounting}
%% ===========================================================

This section derives the accounting law that governs the relationship between
gross accumulation, profitability, capacity utilization, and capital
productivity. Every step is an identity manipulation --- no behavioral equation
is imposed. The behavioral content enters only through the question of whether
the recapitalization rate is constant, a question deferred entirely
to~\S\ref{sec:stageB}.

Define the gross accumulation rate as
\begin{equation}
  \iotat \equiv \frac{I_t}{K_t}.
  \label{eq:iota}
\end{equation}
The accounting decomposition of accumulation is:
\begin{align}
  \iotat &= \phit \cdot \mut \cdot \Bt,
  \qquad \phit \equiv \frac{I_t}{Y_t},
  \label{eq:iota_decomp}
\end{align}
and since $\phit \cdot \mut = (I_t/Y_t)(Y_t/\Yp) = I_t/\Yp \equiv \phi^p_t$,
we have the equivalent form
\begin{equation}
  \iotat = \phi^p_t \cdot \Bt,
  \qquad
  \phi^p_t \equiv \frac{I_t}{\Yp},
  \label{eq:iota_phip}
\end{equation}
where utilization drops out at this level of abstraction.
Introducing the recapitalization rate $\betat \equiv I_t/\Pi_t$,
and using $\pit \equiv \Pi_t/Y_t$ and $r_t \equiv \Pi_t/K_t = \pit\mut\Bt$:
\begin{equation}
  \boxed{\iotat = \betat \cdot \pit \cdot \mut \cdot \Bt = \betat \cdot r_t.}
  \label{eq:iota_full}
\end{equation}

Equation~\eqref{eq:iota_full} is an identity. There are two recapitalization lenses:
\begin{itemize}
  \item $\betat = I_t/\Pi_t$: funding of investment out of profits;
  \item $\phi^p_t = I_t/\Yp$: claim of investment on productive-capacity output.
\end{itemize}
They are linked accounting representations of the same accumulation law.

Both lenses are accounting objects; neither requires a behavioral assumption.
Their analytical value lies in what happens when one asks whether they are
constant. If $\betat$ is constant, accumulation is fully determined by the
accounting decomposition --- capitalists mechanically reinvest a fixed share of
profits regardless of structural conditions. If $\betat$ varies, then the
question becomes: what does it respond to? That is the behavioral margin. Note
that the Weisskopf decomposition~\eqref{eq:weisskopf} is the profitability
reading of the same identity: the profit rate is demand-led through~$\mut$,
structurally mediated through~$\Bt$, and distributively conditioned
through~$\pit$. Section~\ref{sec:stageB} takes up the behavioral margin
directly, specifying the accumulation law in which $\betat$ responds to
distribution, demand, and the structural transformation of capital
productivity.


%% ===========================================================
\section{Stage B: Behavioral Accumulation Law (ARDL)}
\label{sec:stageB}
%% ===========================================================

Stage~B treats the recovered structural objects $(\widehat{\mut}, \hat{B}_t)$
as given conditioning variables and estimates how capitalists behaviorally
determine the recapitalization rate. The behavioral accumulation law is a
conditional module nested inside the structural-identification architecture: it
describes how realized accumulation responds to profitability once the
identification problem has been solved upstream in Stage~A.


%%------------------------------------------------------------
\subsection{Specification}
\label{sec:stageB_spec}
%%------------------------------------------------------------

If recapitalization responds behaviorally to distribution, demand, and
structural conditions, then
\begin{equation}
  \ln\betat = \eta_0 + \eta_1\ln\pit + \eta_2\ln\mut + \eta_3\ln\Bt,
  \label{eq:beta_behavioral}
\end{equation}
which, substituted into the log of~\eqref{eq:iota_full}, yields the
long-run behavioral accumulation law:
\begin{equation}
  \boxed{
  \ln\iotat = c + b_1\ln\pit + b_2\ln\mut + b_3\ln\Bt + \varepsilon_t
  }
  \label{eq:accumulation_law}
\end{equation}
with $b_j = 1 + \eta_j$.

\begin{proposition}[Dual Interpretation of Coefficients]
\label{prop:dual}
Each coefficient in~\eqref{eq:accumulation_law} has both an accounting
and a behavioral reading:
\end{proposition}

\begin{table}[h]
  \centering
  \caption{Dual Interpretation of Accumulation Law Coefficients}
  \label{tab:dual}
  \begin{tabular}{lll}
    \toprule
    Coefficient & Accounting reading & Behavioral reading \\
    \midrule
    $b_1$ & Profit-share elasticity of accumulation
           & $\eta_1 = b_1 - 1$: recapitalization response to distribution \\
    $b_2$ & Utilization elasticity of accumulation
           & $\eta_2 = b_2 - 1$: recapitalization response to demand \\
    $b_3$ & Capital-productivity elasticity of accumulation
           & $\eta_3 = b_3 - 1$: recapitalization response to structural conditions \\
    \bottomrule
  \end{tabular}
  \bigskip

  \noindent{\small \textit{Notes}: Each $b_j = 1 + \eta_j$.
  The deviation from unity is the behavioral content.
  Estimation uses ARDL with Pesaran--Shin--Smith (2001) bounds test.}
\end{table}


%%------------------------------------------------------------
\subsection{Nested Restriction Menu}
\label{sec:restrictions}
%%------------------------------------------------------------

The behavioral coefficients generate a hierarchy of accumulation closures:

\begin{table}[h]
  \centering
  \caption{Nested Restriction Menu}
  \label{tab:restrictions}
  \begin{tabular}{lll}
    \toprule
    Restriction & Null & Model name \\
    \midrule
    Cambridge closure       & $b_1 = b_2 = b_3 = 1$ & $\betat$ constant \\
    Bhaduri--Marglin type   & $b_3 = 1$; $b_1, b_2$ free & Blind to capital productivity \\
    Full structural model   & All $b_j$ free & All channels active \\
    \bottomrule
  \end{tabular}
\end{table}


%%------------------------------------------------------------
\subsection{Results: US}
\label{sec:stageB_us}
%%------------------------------------------------------------

\draftnote{§2.8.3: 2--3 pages (to be filled after estimation).
  ARDL estimation with PSS bounds test. Report $\hat{b}_1, \hat{b}_2, \hat{b}_3$
  and associated Wald tests. Which closure is supported?
  Interpret $\eta_j = b_j - 1$ economically.}


%%------------------------------------------------------------
\subsection{Results: Chile}
\label{sec:stageB_chile}
%%------------------------------------------------------------

\draftnote{§2.8.4: 2--3 pages (to be filled after estimation).
  Same procedure. Key question: does $b_3$ behave differently from the US?
  Does the capital-productivity channel operate freely or is it structurally
  constrained? Transition: §2.9 provides the mechanism.}


%% ===========================================================
\section{The Peripheral Twist: BoP Constraint and the Import
  Content of Recapitalization}
\label{sec:tavares}
%% ===========================================================

This section develops the chapter's distinctive theoretical contribution. The
structural comparison of the behavioral accumulation law rests on a specific
claim about why the recapitalization response differs between center and
periphery. That claim originates in the work of
\citeauthor{tavares1985} (\citeyear{tavares1985}, \citeyear{tavares2000}): the
binding constraint on peripheral accumulation is not on output directly but on
the ability to convert surplus into re-investment, because the capital goods
required for accumulation must be imported.


%%------------------------------------------------------------
\subsection{The Tavares Channel}
\label{sec:tavares_channel}
%%------------------------------------------------------------

The first step of the mechanism is structural incompleteness of the
capital-goods sector. Peripheral industrialization, whether under
import-substitution regimes or their liberalized successors, is characterized
by what \citet{tavares1985} calls the ``hard phase'' of industrialization and
what \citet{fajnzylber1983} terms ``truncated industrialization'': the domestic
economy does not produce the full range of machinery and equipment required for
capital formation. The share of machinery and equipment in gross capital
formation that must be sourced from abroad ---
denoted~$\tilde{\phi}^{ME}_t$ --- is not a policy variable that governments
can adjust at will but a structural parameter reflecting the depth and
composition of domestic manufacturing capacity. In Latin America's larger
economies, this share has oscillated historically between 0.30 and 0.60
depending on the phase of the industrialization cycle, but it has never
approached zero. Even during the most ambitious import-substitution episodes,
the capital-goods sector remained incomplete in the sense that the most
technologically advanced vintages of equipment continued to require foreign
sourcing. The consequence is that gross investment has an irreducible imported
component, and the domestic availability of profits --- however large --- cannot
by itself finance the physical act of capital formation.

The second step links this structural incompleteness to the recapitalization
rate. The accounting identity $\iotat = \betat\, r_t$ established
in~\S\ref{sec:accounting} decomposes the gross accumulation rate into the
recapitalization rate and the profit rate. In the center, $\betat$ responds to
the structural conditions of accumulation --- to distribution, demand, and
capital productivity --- with relatively few binding external constraints on the
conversion of profits into investment goods. In the periphery, however, the
conversion channel is mediated by the foreign-exchange market. Profits exist
($\Pi_t > 0$), the profit rate may be high ($r_t$ may exceed the center's),
and the desire to reinvest may be strong --- yet $\betat$ collapses when the
foreign currency required to purchase imported machinery is unavailable or
prohibitively expensive. The surplus exists but cannot be converted into
re-investment. This is the core of Tavares's insight, articulated most clearly
in the retrospective synthesis of \citet{tavares2000} as read through
\citet{vernengo2006}: the binding constraint on peripheral accumulation is not
the generation of surplus but its physical realization as installed productive
capital.

The third step identifies the historical transmission mechanism. The constraint
on $\betat$ does not operate continuously but activates at discrete episodes of
external financial stress --- debt-cycle reversals, terms-of-trade collapses,
sudden stops in capital inflows. The Volcker shock of 1979--1982 is the
canonical example for Latin America: the sudden repricing of dollar-denominated
debt simultaneously drained foreign-exchange reserves and raised the
domestic-currency cost of imported capital goods, producing a collapse in
$\betat$ that persisted for nearly a decade despite the partial recovery of
domestic profitability. These episodes express financial dependency as a
constraint on the technical composition of accumulated capital. They are not
generic recessions in which output and investment fall together; they are
structurally asymmetric crises in which the capacity to generate surplus and the
capacity to convert that surplus into productive capital are severed by the
foreign-exchange constraint.

It is essential to distinguish the Tavares channel from two superficially
similar but analytically distinct mechanisms. It is not a pure output-side
balance-of-payments constraint of the \citet{thirlwall1979} variety, in which
export growth determines the growth of output through the income elasticity of
imports --- although the two may operate simultaneously. And it is not simple
import compression, in which a foreign-exchange shortage reduces the volume of
all imports proportionally. The Tavares channel is specifically a
recapitalization constraint: it operates on $\betat$, not on $Y_t$, and it
binds through the import content of investment goods, not through the aggregate
import bill. The contribution of this chapter is to give that theoretical
distinction empirical content by estimating the behavioral accumulation law
separately for the US and Chile and testing whether the coefficient $b_3$ ---
the recapitalization response to capital productivity --- is structurally
constrained in the periphery relative to the center.


%%------------------------------------------------------------
\subsection{Formalization}
\label{sec:tavares_formal}
%%------------------------------------------------------------

The Tavares channel described in~\S\ref{sec:tavares_channel} identifies a
structural constraint on $\betat$ in the periphery, operating through the
import content of investment and the availability of foreign exchange. The key
claim is that the recapitalization response to capital productivity --- the
parameter $\eta_3 = b_3 - 1$ in the behavioral accumulation
law~\eqref{eq:accumulation_law} --- is not a free behavioral parameter in the
periphery. It is structurally constrained by two forces: the import content of
investment~$\tilde{\phi}^{ME}_t$ and the external financial conditions
summarized by a balance-of-payments proxy $\text{BoP}_t$.

Disaggregated import data for Chile at the machinery-and-equipment level are
not available at annual frequency over the full sample, so the empirical
strategy adopts the \citet{kaldor1959} prior: from the cross-country
tabulation in Cuadro~8, the machinery-and-equipment share of capital formation
in semi-industrialized economies clusters around
$\xi_K^{ME} \approx 0.92$--$0.94$. This value is imposed as an empirical
prior, not estimated from unavailable disaggregated data. The balance-of-payments
proxy~$\text{BoP}_t$ admits several candidate series --- the terms of trade,
the real exchange rate, or the current-account balance --- and the choice among
them is an empirical question adjudicated by the Stage~B estimation.

Three candidate specifications implement the constraint. The first augments
the Chilean ARDL with a foreign-exchange availability proxy, entering
$\text{BoP}_t$ directly as an additional regressor in the behavioral
accumulation law and testing whether it absorbs part of the $b_3$ channel. The
second tests whether $b_3^{CL}$ is stable across balance-of-payments regimes
or breaks at sudden-stop episodes --- the preferred implementation is the
\citet{hansenseo2002} threshold VECM, in which the balance-of-payments
constraint activates as a threshold on $\tilde{\phi}^{ME}_t$ rather than
operating linearly throughout the sample. The third compares the unconstrained
US estimate~$b_3^{US}$ directly against the Chilean estimate~$b_3^{CL}$ and
tests the structural null $\eta_3^{US} > \eta_3^{CL}$: the peripheral
recapitalization response to capital productivity is more constrained than the
center's.

It is important to distinguish two sources of external constraint that the
literature sometimes conflates. The income elasticity of exports~$\varepsilon$
is an exogenous parameter reflecting the composition of global trade. The
import propensity is endogenous to the domestic productive structure: it
reflects the degree of structural incompleteness of the capital-goods sector.
The Kaldor-versus-ECLA fault line turns on whether the external ceiling on
accumulation is determined primarily by $\varepsilon$ (the export side, as in
Thirlwall) or by the import propensity operating through
$\tilde{\phi}^{ME}_t$ (the import-content side, as in Tavares). This is an
estimable structural question, and the threshold-VECM specification is designed
to adjudicate it.

Let $\tilde{\phi}^{ME}_t$ denote the import content of investment —
the share of machinery and equipment in gross capital formation that is
sourced from abroad. The peripheral behavioral law gains a
forex-mediation channel:
\begin{equation}
  \eta^{CL}_3 = \eta^{CL}_3\!\left(\tilde{\phi}^{ME}_t, \text{BoP}_t\right),
  \label{eq:tavares_constraint}
\end{equation}
where $\text{BoP}_t$ is a measure of external constraint intensity
(terms of trade, real exchange rate, or current account balance).


%%------------------------------------------------------------
\subsection{The Technical Composition Channel}
\label{sec:tech_composition}
%%------------------------------------------------------------

The Tavares channel operates through the recapitalization rate; the technical
composition channel explains why over-mechanization amplifies the constraint.
When the capacity transformation elasticity satisfies $\theta < 1$ --- the
over-mechanization regime identified in~\S\ref{sec:primitives} --- two
consequences follow simultaneously, and their interaction produces a structural
trap specific to the periphery.

The first consequence is that more capital is required per unit of productive
capacity created. Under over-mechanization, the capital stock grows faster than
productive capacity: $\ln\Bt = (\theta - 1)k_t < 0$ when $\theta < 1$, so
capital productivity at normal capacity falls secularly. Each successive unit
of capacity requires a larger capital outlay. In an economy with a complete
domestic capital-goods sector, this is a drag on accumulation but not a binding
constraint --- the additional capital goods can be produced domestically,
financed from the surplus, and installed without crossing the foreign-exchange
market. In the periphery, however, the structural incompleteness described
in~\S\ref{sec:tavares_channel} means that the additional capital goods must be
imported. Over-mechanization does not merely raise the volume of investment
required; it raises the volume of \emph{imported} investment required, because
the marginal capital goods needed to sustain capacity growth are
disproportionately those that the domestic economy does not produce.

The second consequence is that the economy's capacity to earn the foreign
exchange needed to finance those imports deteriorates. If capital productivity
falls, the output obtainable from the installed capital stock at normal
utilization declines, and with it the export base that generates
foreign-currency earnings. The secular deterioration of the terms of trade
identified by Prebisch and Singer is not, in this framework, merely a price
phenomenon --- it is a structural constraint on the technical composition of
peripheral capital accumulation. The economy needs more imported machinery per
unit of capacity, but its capacity to earn the forex that pays for that
machinery is lower because the accumulated capital is less productive.

The peripheral trap is the interaction of these two forces.
Over-mechanization raises the import bill for investment while degrading the
export capacity that finances it. The $b_3$ comparison
in~\S\ref{sec:comparison} gives this trap empirical content: if
$b_3^{CL} < b_3^{US}$, then Chilean capitalists cannot recapitalize at the
rate the structural conditions demand, because the foreign-exchange constraint
binds. The finding would indicate that the recapitalization response to capital
productivity is structurally dampened in the periphery --- not because Chilean
capitalists lack the will or the surplus to reinvest, but because the physical
conversion of surplus into installed productive capital requires imported
machinery that the balance of payments cannot finance.

This is the chapter's central structural claim about the center--periphery
difference: \emph{demand conditions intervene on the political economy of the
ceiling, i.e.\ the class struggle over distribution, utilization, and the
institutional forms that govern whether accumulation can proceed.}


%% ===========================================================
\section{Structural Comparison: Center vs.\ Periphery}
\label{sec:comparison}
%% ===========================================================

\draftnote{§2.10 intro: 1 paragraph. Synthesize the results across both countries.
  The parallel construction (§2.5.3) ensures the comparison is structural,
  not an artefact of measurement.}

\begin{table}[h]
  \centering
  \caption{Center--Periphery Structural Comparison}
  \label{tab:comparison}
  \begin{tabular}{llll}
    \toprule
    Diagnostic & US & Chile & Reading \\
    \midrule
    $\hat{\theta}$               & --- & --- & Mechanization gap \\
    Regime breaks                & --- & --- & When does $\theta$ shift? \\
    $b_1$ (distribution)         & --- & --- & Profit-share response \\
    $b_2$ (demand)               & --- & --- & Utilization response \\
    $b_3$ (structural)           & --- & --- & Capital-productivity response \\
    $b_3$ stability              & --- & --- & Does $b_3^{CL}$ break at BoP crises? \\
    $\widehat{\mu}_t$ volatility & --- & --- & Utilization instability \\
    $\hat{B}_t$ trend            & --- & --- & Capital productivity trajectory \\
    \bottomrule
  \end{tabular}
  \bigskip

  \noindent{\small \textit{Notes}: Cells filled after Stage A and Stage B
  estimation. $b_3$ is the key comparison: in the US, the recapitalization
  function responds to the capital-productivity channel relatively freely;
  in Chile, $b_3^{CL}$ is mediated by the BoP constraint (§2.9).}
\end{table}

\draftnote{§2.10 narrative: explain each row of Table~\ref{tab:comparison}.
  Punchline: financial dependency shows up empirically as a structurally
  constrained recapitalization response. The Tavares channel provides the
  mechanism; Table~\ref{tab:comparison} provides the evidence.}


%% ===========================================================
\section{Conclusions}
\label{sec:conclusions}
%% ===========================================================

\draftnote{§2.11: approx.\ 3 pages. Three deliverables:
  (1)~CU series: structurally identified $\hat{\mu}_t$ for both countries,
  available as inputs to Ch3.
  (2)~Regime classification: $\hat{\theta}^{US}$ vs.\ $\hat{\theta}^{CL}$
  — the empirical realization of the core-periphery mechanization gap.
  (3)~Financial dependency formalized: the Tavares channel — inability to
  convert surplus into re-investment due to lack of foreign currency —
  is not a narrative but an estimable structural constraint on the
  recapitalization function.
  Forward references: Ch3 inherits $\hat{\mu}_t$, $\hat{B}_t$, $\hat{\theta}_t$
  as given inputs for profitability analysis and investment function
  hypothesis testing. The Weisskopf decomposition and Okishio crisis
  trigger become Ch3 applications of the Stage A recovered objects.}


%% ===========================================================
%%  BIBLIOGRAPHY
%% ===========================================================
\bibliographystyle{plainnat}
\bibliography{ch2_refs}

%% Placeholder references (to be migrated to .bib file)
%
% \bibitem Okishio (1961): "Technical Change and the Rate of Profit"
% \bibitem Weisskopf (1979): "Marxian crisis theory and the rate of profit"
% \bibitem Tavares (1985): "Auge e declínio do processo de substituição"
% \bibitem Vernengo (2006): "Technology, finance, and dependency"
% \bibitem Basu & Das (2017): "Foreign capital, fiscal policy, and crisis"
% \bibitem Vidal (2014, 2019): "Stagnation tendencies"
% \bibitem Basu (2019): "Financialization, stagnation"
% \bibitem Hansen-Seo (2002): threshold VECM
% \bibitem Kaldor (1959): "Problemas Económicos de Chile"
% \bibitem Palma & Marcel (1989): surplus and development

\end{document}
